% Table rows are generated from recorded runs by Sources/scripts/journal_experiments/fill_manuscript.py.
% \tabrows uses the TeX primitive input so it works inside tabular.
% \figorph{file}{text}: include the generated figure if it exists, else keep the placeholder box.
\providecommand{\figorph}[2]{\IfFileExists{#1}{\includegraphics[width=\linewidth]{#1}}{\fbox{\parbox[c][5cm][c]{0.9\linewidth}{\centering #2}}}}
\makeatletter\providecommand{\tabrows}[1]{\@@input #1 }\makeatother
\InputIfFileExists{generated/prose_values}{}{}
\section{Experimental Evaluation}
\label{sec:experiments}

\subsection{Experimental Setup}
The experimental campaign evaluates \ehsfp under hierarchical split-federated learning with complete, partial, stale, and episodic prototype arrivals. Journal-protocol classification experiments use CIFAR-10 and CIFAR-100~\cite{krizhevsky2009learning} and HAM10000~\cite{tschandl2018ham10000}; ImageNet-1K~\cite{5206848} appears only in the conference-reference tables, because no journal-protocol ImageNet-1K runs were performed. Segmentation experiments use ISIC-2018, following the skin-lesion challenge protocol lineage in~\cite{8363547}. The compared methods are FedAvg~\cite{mcmahan2017communication}, FedProx~\cite{li2020federated}, FedNova~\cite{10.5555/3495724.3496362}, HierFL~\cite{liu2020client,9488756}, SplitFed~\cite{thapa2022splitfed}, HeteroSFL, HSFL~\cite{quan2023hiersfl}, the conference \hsfp baseline summarized in Section~\ref{sec:preliminary}, and \ehsfp.

Data are partitioned into IID, mildly non-IID, and severely non-IID client splits. Non-IID splits use class-proportion Dirichlet sampling with concentration $0.7$ (mild) and $0.3$ (severe), fixed random seeds, and identical client counts across compared methods. % TODO(result): Confirm the exact Dirichlet concentration values for mild and severe non-IID partitions after the simulation configuration is frozen.
The episodic stress protocol varies prototype dropout rate $\dropoutp$, staleness bound $\stale$, asynchronous edge/cloud aggregation interval, and serverless-style event loss, timeout, and cold-start conditions. These stressors are included to test whether stateless learning episodes can remain stable when prototype packets arrive incompletely or late.

Evaluation metrics include accuracy, macro F1-score, communication volume, client memory, edge memory, cloud memory, rounds to convergence, training time, Dice, IoU, prototype drift, and stability. Prototype drift follows Definition~\ref{def:proto_drift}. For a validation-accuracy sequence $\{a_t\}$ and a trailing window $\mathcal{W}_t$, the stability metric is defined as
\begin{equation}
    \mathrm{Stability}(t)
    =
    \frac{1}{|\mathcal{W}_t|}
    \sum_{s\in\mathcal{W}_t}
    \left(a_s-\frac{1}{|\mathcal{W}_t|}\sum_{r\in\mathcal{W}_t}a_r\right)^2 .
\end{equation}
Lower values indicate smaller validation oscillation over the trailing window. We use $|\mathcal{W}_t|=5$ and report the value at the final round in units of $\%^2$. Rounds to convergence is the first round at which the three-round moving average of validation accuracy reaches 95\% of the best accuracy observed in the run and remains there for three rounds. Prototype drift is the edge-tier squared 2-Wasserstein drift of Definition~\ref{def:proto_drift}, averaged over edges and rounds. Communication is the total traffic over the run in GB.

\paragraph{Implementation.}
All journal-protocol classification runs use 200 clients grouped under 5 edge servers, 10\% client participation per round, 60 communication rounds, a local batch size of 16, and Adam with learning rate $10^{-4}$, weight decay $10^{-4}$, and a five-round warm-up followed by a learning-rate schedule. On CIFAR-100, the client holds an 11.2M-parameter ResNet-18 trunk, the edge holds a 0.69M-parameter module, and the cloud holds a linear classifier. The full \ehsfp configuration uses episodic prototype memory (capacity 500, replay ratio 0.3, maximum prototype age 20), prototype dropout with rate 0.2, reliability-aware aggregation, replay consistency, and low-rank shared-covariance synthesis. It also runs with simulated serverless events: cold starts with probability 0.15 and function timeouts with probability 0.05. The \hsfp rows use the same backbone, partitions, and schedule but no simulated events. Consequently, \hsfp--\ehsfp accuracy comparisons are conservative for \ehsfp, while their prototype-drift values are not directly comparable. % TODO(result): Confirm hardware and the edge/cloud intervals before submission.

\paragraph{Protocol separation and result provenance.}
We keep three result streams strictly disjoint. First, the \emph{conference-reference protocol} contains the five-run results reported in the accepted ECCV 2026 paper, including 200 distributed devices and the aggregation intervals stated in each table caption. Second, the \emph{proxy protocol} uses 10 communication rounds and three date-based seeds (20260714--20260716); it is used only for implementation diagnosis and is not used for paper-level comparisons. Third, the \emph{frozen journal protocol} uses 60 rounds and integer seeds. The matched method comparisons (Tables~\ref{tab:journal_main_classification} and~\ref{tab:journal_segmentation}) use five seeds, $\{0,1,2,3,4\}$. The robustness, ablation, and scalability sweeps use a single run with seed 0 per cell. The convergence curves in Figure~\ref{fig:conv_curves} use seeds $\{0,1,2\}$. Single-seed differences smaller than the five-seed standard deviations of Table~\ref{tab:journal_main_classification} (0.3--0.6 points on CIFAR) should therefore not be read as effects. Consequently, proxy, legacy 60-round, and conference-reference values are never pooled, rescaled, or compared as if they were generated by a common protocol. Unless otherwise stated, reported uncertainty is the sample standard deviation over independent runs.

% INTERNAL STATUS (2026-07-22): no completed Full E-HSFP result is available.
% The hsfp_memory_reliability_prc and full_e_hsfp presets are blocked by an
% unresolved autograd deadlock. Keep E-HSFP cells as \pending until genuine,
% traceable five-seed results exist under the frozen journal protocol.

\subsection{Convergence Behavior}
This subsection connects the empirical campaign to Theorem~\ref{thm:convergence} and Corollary~\ref{cor:small_perturb}. Figure~\ref{fig:conv_curves} shows validation curves under IID and non-IID partitions (mean over seeds $\{0,1,2\}$). At round 60, \ehsfp is ahead of \hsfp on CIFAR-10 IID (49.25 vs.\ 45.61\%) and $dir(0.7)$ (39.70 vs.\ 36.61\%). Under $dir(0.3)$ the two are indistinguishable (35.30$\pm$3.24 vs.\ 35.22$\pm$0.14\%), with a much larger seed-to-seed spread for \ehsfp. On CIFAR-100 IID, \ehsfp leads 48.24 vs.\ 44.96\%. On HAM10000, \ehsfp reaches 67.27--70.66\% across the three partitions, whereas \hsfp stays at 40.44--54.40\% with standard deviations of up to 7.6 points. The seed-0 \hsfp runs on HAM10000 never meet the convergence rule above (Table~\ref{tab:noniid_robustness}). Figure~\ref{fig:drift_rounds} shows prototype drift over rounds under the episodic stress protocol. The stability-bound term of Figure~\ref{fig:bound_vs_drift} has not been evaluated numerically, because the constants of Theorem~\ref{thm:convergence} have not been instantiated for these models; the corresponding column of Table~\ref{tab:staleness_convergence} is marked n/a.

% TODO(figure): Replace placeholder with final plot after simulations are completed.
\begin{figure}[t]
    \centering
    \figorph{figures/fig_conv_curves.pdf}{Placeholder: convergence curves under IID and non-IID partitions (final plot pending simulations)}
    \caption{Validation convergence curves for IID and non-IID classification protocols.}
    \label{fig:conv_curves}
\end{figure}

% TODO(figure): Replace placeholder with final plot after simulations are completed.
\begin{figure}[t]
    \centering
    \figorph{figures/fig_drift_rounds.pdf}{Placeholder: prototype drift over communication rounds (final plot pending simulations)}
    \caption{Prototype-induced distributional drift over rounds under the episodic stress protocol.}
    \label{fig:drift_rounds}
\end{figure}

% TODO(figure): Replace placeholder with final plot after simulations are completed.
\begin{figure}[t]
    \centering
    \fbox{\parbox[c][5cm][c]{0.9\linewidth}{\centering Placeholder: stability-bound term versus empirical drift (final plot pending simulations)}}
    \caption{Comparison between the stability-bound term from the convergence analysis and empirical prototype drift.}
    \label{fig:bound_vs_drift}
\end{figure}

\subsection{Main Results}
Tables~\ref{tab:main_classification} and~\ref{tab:segmentation} reproduce the accepted conference-reference results verbatim; they establish the H-SFP baseline that the journal extension must compare against. These values are not pooled with the proxy or legacy 60-round diagnostics. Under the matched conference protocol, H-SFP provides the highest federated accuracy on all four classification datasets. It exceeds the strongest reported federated baseline by 3.99 percentage points on CIFAR-10, 2.00 points on CIFAR-100, 7.76 points on HAM10000, and 1.1 points on ImageNet-1K.

Table~\ref{tab:journal_main_classification} repeats the comparison under the frozen journal protocol with five seeds. \ehsfp improves on \hsfp on every dataset: by 3.71 points on CIFAR-10 (49.14$\pm$0.27 vs.\ 45.43$\pm$0.27), 3.19 points on CIFAR-100 (48.05$\pm$0.59 vs.\ 44.86$\pm$0.30), and 25.80 points on HAM10000 (67.21$\pm$1.15 vs.\ 41.41$\pm$2.74). This improvement holds even though the \ehsfp runs include simulated serverless events. However, the conference ranking is not reproduced under this protocol. Both prototype-based methods trail the strongest weight-sharing baselines: HSFL (A) reaches 73.75\% on CIFAR-10 and 58.24\% on CIFAR-100, and SplitFed reaches 78.31\% on HAM10000. FedAvg, FedProx, HierFL, and SplitFed also exceed \ehsfp on all three datasets. \ehsfp is ahead only of FedNova and HeteroSFL, both of which degrade severely under this protocol. The journal-protocol \hsfp is also far below its conference-reference accuracy (45.43 vs.\ 67.44\% on CIFAR-10), which suggests that the gap originates in the shared prototype pipeline under this backbone and budget rather than in the E-HSFP modules. The present evidence therefore supports \ehsfp as an improvement over \hsfp under episodic stress, but not as a state-of-the-art federated classifier.

\begin{table*}[t]
\centering
\caption{Conference-reference classification accuracy (\%) across 200 distributed devices. CIFAR-10 and CIFAR-100 use IID partitions; HAM10000 and ImageNet-1K use ResNet-50 under IID partitions. HSFL and H-SFP use aggregation interval (A), $I_c=5$ and $I_e=10$. Results are mean $\pm$ standard deviation over five independent runs and are reproduced verbatim from the accepted ECCV 2026 paper. E-HSFP remains pending under the frozen journal protocol.}
\label{tab:main_classification}
\begin{tabular}{@{}lcccc@{}}
\toprule
Method & CIFAR-10 & CIFAR-100 & HAM10000 & ImageNet-1K \\
\midrule
FedAvg    & 57.25 {\scriptsize $\pm$ 1.24} & 48.85 {\scriptsize $\pm$ 1.40} & 68.28 {\scriptsize $\pm$ 1.45} & 23.1 {\scriptsize $\pm$ 1.4} \\
FedProx   & 58.25 {\scriptsize $\pm$ 1.20} & 52.75 {\scriptsize $\pm$ 1.25} & 71.15 {\scriptsize $\pm$ 1.30} & 24.1 {\scriptsize $\pm$ 1.3} \\
FedNova   & 23.50 {\scriptsize $\pm$ 1.15} & 50.24 {\scriptsize $\pm$ 1.35} & 69.60 {\scriptsize $\pm$ 1.35} & 23.8 {\scriptsize $\pm$ 1.3} \\
FedProto  & 55.10 {\scriptsize $\pm$ 1.05} & 46.50 {\scriptsize $\pm$ 1.15} & 67.20 {\scriptsize $\pm$ 1.25} & 25.2 {\scriptsize $\pm$ 1.2} \\
HierFL    & 62.50 {\scriptsize $\pm$ 1.30} & 30.15 {\scriptsize $\pm$ 1.45} & 69.50 {\scriptsize $\pm$ 1.25} & 24.5 {\scriptsize $\pm$ 1.2} \\
SplitFed  & 61.85 {\scriptsize $\pm$ 1.25} & 35.00 {\scriptsize $\pm$ 1.30} & 65.78 {\scriptsize $\pm$ 1.35} & 25.5 {\scriptsize $\pm$ 1.1} \\
HeteroSFL & 63.45 {\scriptsize $\pm$ 1.10} & 41.50 {\scriptsize $\pm$ 1.25} & 67.10 {\scriptsize $\pm$ 1.20} & 26.3 {\scriptsize $\pm$ 0.8} \\
FedGen    & 59.80 {\scriptsize $\pm$ 1.10} & 51.20 {\scriptsize $\pm$ 1.25} & 72.50 {\scriptsize $\pm$ 1.15} & 26.8 {\scriptsize $\pm$ 1.0} \\
FedDF     & 61.50 {\scriptsize $\pm$ 1.15} & 53.10 {\scriptsize $\pm$ 1.20} & 73.10 {\scriptsize $\pm$ 1.20} & 25.9 {\scriptsize $\pm$ 1.1} \\
HSFL (A)  & 51.35 {\scriptsize $\pm$ 1.10} & 35.47 {\scriptsize $\pm$ 1.35} & 68.82 {\scriptsize $\pm$ 1.25} & 24.6 {\scriptsize $\pm$ 0.9} \\
\textbf{H-SFP (A)} & \textbf{67.44 {\scriptsize $\pm$ 0.85}} & \textbf{55.10 {\scriptsize $\pm$ 0.95}} & \textbf{80.86 {\scriptsize $\pm$ 0.75}} & \textbf{27.9 {\scriptsize $\pm$ 0.3}} \\
E-HSFP & \pending & \pending & \pending & \pending \\
\bottomrule
\end{tabular}
\end{table*}

\begin{table}[t]
\centering
\caption{Journal-protocol matched classification accuracy (\%): 200 clients, 5 edges, IID, 60 rounds, seeds $\{0,\dots,4\}$; final-round (round 60) test accuracy for every method (no test-set model selection). Mean $\pm$ sample standard deviation; a superscript $n$ marks cells with fewer than five completed seeds. These values are not pooled with Table~\ref{tab:main_classification}. $^\ddagger$HierFL selects propagated models using test-split evaluations during training. $^\S$The cell includes seeds whose training collapsed (constant or near-chance prediction); they are reported, not discarded.}
\label{tab:journal_main_classification}
\begin{tabular}{@{}lccc@{}}
\toprule
Method & CIFAR-10 & CIFAR-100 & HAM10000 \\
\midrule
\tabrows{generated/journal_main_classification.tex}
\bottomrule
\end{tabular}
\end{table}

On ISIC-2018, H-SFP improves over HSFL by 6.2 points in IoU and 4.7 points in Dice while reducing training time from 14.5 to 6.8 hours. Relative to SplitFed, H-SFP improves IoU by 14.5 points and is approximately $10.6\times$ faster. These are conference-reference comparisons under a common 200-round, 200-device protocol. Under the journal protocol (Table~\ref{tab:journal_segmentation}), \ehsfp does not improve segmentation. It reaches 45.89$\pm$1.67\% IoU and 61.01$\pm$1.67\% Dice, which is 6.27 and 5.99 points below \hsfp (52.16$\pm$1.64 and 67.00$\pm$1.57), and it needs 8.28 instead of 5.51 hours on the same host. HierFL attains the highest IoU (71.54$\pm$1.50), but its model selection uses test-split evaluations during training.

\begin{table}[t]
\centering
\caption{Conference-reference ISIC-2018 segmentation results using a ResNet50-UNet backbone, 200 communication rounds, and 200 distributed devices. Results are mean $\pm$ standard deviation over five runs. E-HSFP remains pending under the frozen journal protocol.}
\label{tab:segmentation}
\begin{tabular}{@{}lccc@{}}
\toprule
Method & IoU (\%) $\uparrow$ & Dice (\%) $\uparrow$ & Time (h) $\downarrow$ \\
\midrule
FedAvg    & 48.7 {\scriptsize $\pm$ 1.2} & 64.2 {\scriptsize $\pm$ 1.5} & 28.5 {\scriptsize $\pm$ 0.8} \\
HierFL    & 51.2 {\scriptsize $\pm$ 1.1} & 66.3 {\scriptsize $\pm$ 1.3} & 24.3 {\scriptsize $\pm$ 0.7} \\
SplitFed  & 53.8 {\scriptsize $\pm$ 0.9} & 68.1 {\scriptsize $\pm$ 1.1} & 72.4 {\scriptsize $\pm$ 1.6} \\
HeteroSFL & 56.4 {\scriptsize $\pm$ 0.7} & 70.5 {\scriptsize $\pm$ 0.8} & 45.8 {\scriptsize $\pm$ 1.2} \\
HSFL (A)  & 62.1 {\scriptsize $\pm$ 0.8} & 74.8 {\scriptsize $\pm$ 0.9} & 14.5 {\scriptsize $\pm$ 0.5} \\
\textbf{H-SFP (A)} & \textbf{68.3 {\scriptsize $\pm$ 0.5}} & \textbf{79.5 {\scriptsize $\pm$ 0.4}} & \textbf{6.8 {\scriptsize $\pm$ 0.2}} \\
E-HSFP & \pending & \pending & \pending \\
\bottomrule
\end{tabular}
\end{table}

\begin{table}[t]
\centering
\caption{Journal-protocol ISIC-2018 segmentation (50 clients, 60 rounds, seeds $\{0,\dots,4\}$); each runner's reported test IoU/Dice and host wall-clock time. Not pooled with Table~\ref{tab:segmentation}. $^\ddagger$See Table~\ref{tab:journal_main_classification}. $^\star$Wall-clock time measured on the previous host (RTX 3080 Ti, up to four concurrent runs); not comparable with times measured on the current host (RTX 6000 Ada, twelve concurrent runs).}
\label{tab:journal_segmentation}
\begin{tabular}{@{}lccc@{}}
\toprule
Method & IoU (\%) $\uparrow$ & Dice (\%) $\uparrow$ & Time (h) $\downarrow$ \\
\midrule
\tabrows{generated/journal_segmentation.tex}
\bottomrule
\end{tabular}
\end{table}

% INTERNAL DIAGNOSTIC NOTE (not a manuscript result): the 10-round ISIC proxy
% grid gives IoU 63.17+-1.56 (HierFL), 47.27+-2.19 (H-SFP), 41.44+-0.52
% (SplitFL), 34.50+-4.33 (HSFL), 27.76+-0.86 (Federated), and
% 25.80+-0.26 (HeteroSFL). Because it uses a different round budget and seed
% convention, it must not replace or be averaged with Table~\ref{tab:segmentation}.

\subsection{Robustness under Episodic Updates}
This subsection evaluates robustness when upper-tier episodes receive partial, stale, or disrupted prototype packets. All cells are single runs with seed 0 on CIFAR-100 IID unless stated otherwise, and the reference configuration is the full \ehsfp model (48.52\% accuracy).

\paragraph{Prototype dropout.} Table~\ref{tab:dropout_convergence} and Figure~\ref{fig:dropout_effect} show that final accuracy is essentially insensitive to the prototype dropout rate. Across $p\in\{0,0.1,0.2,0.3,0.5\}$, accuracy stays within 48.40--48.88\%, a range smaller than the five-seed standard deviation. Prototype drift increases monotonically with $p$, from 167.7 to 171.6, while rounds to convergence vary between 37 and 42.

\paragraph{Stale prototypes.} Staleness has a clearer effect (Table~\ref{tab:staleness_convergence}, Figure~\ref{fig:staleness_effect}). Any staleness $\tau\geq1$ lowers accuracy by 1.5--2.0 points relative to $\tau=0$ (46.52--47.03 vs.\ 48.52\%). Beyond $\tau=1$, however, the loss does not grow further. Prototype drift rises monotonically with $\tau$, from 169.2 at $\tau=0$ to 197.8 at $\tau=10$, which is consistent with the stale-update term in Theorem~\ref{thm:convergence}.

\paragraph{Serverless-style disruptions.} Relative to the run without event loss (48.76\%), function timeouts, partial edge execution, and missing edge updates cost at most 0.42 points (Table~\ref{tab:serverless_robustness}). Cold-start delay costs 1.81 points and the combined stress 1.50 points. The combined stress also produces the largest prototype drift (180.6 vs.\ 159.6 without event loss).

\paragraph{Statistical heterogeneity.} Table~\ref{tab:noniid_robustness} compares \hsfp and \ehsfp across partitions. On CIFAR-10, \ehsfp is ahead under IID ($+3.85$ points) and $dir(0.7)$ ($+3.72$). Under $dir(0.3)$ the seed-0 run is behind ($-3.79$), but the three-seed means in Figure~\ref{fig:conv_curves} are equal, so this cell reflects seed variance rather than a systematic loss. On HAM10000, \ehsfp improves accuracy by 21.0--27.5 points and macro-F1 by 11.5--15.2 points in every partition. \hsfp on CIFAR-100 is insensitive to the partition (44.62--45.14\%). On CIFAR-100, \ehsfp reaches \PVcHundredMildEhsfp{} under $dir(0.7)$ and \PVcHundredSevereEhsfp{} under $dir(0.3)$, compared with 44.82 and 45.14\% for \hsfp and 48.52\% for \ehsfp under IID. In every partition, \ehsfp uses slightly less communication than \hsfp, for example 120.1 vs.\ 125.5~GB on CIFAR-100 IID. It reports higher prototype drift, which reflects the simulated packet loss and staleness included in its runs and should not be read as a matched comparison.

% TODO(result): Fill this table after simulations are completed. Do not fabricate numbers.
\begin{table*}[t]
\centering
\caption{Prototype-dropout robustness under the episodic update protocol. Single run (seed 0) per cell, including reference rows; journal protocol, 60 rounds, final-round test metrics.}
\label{tab:dropout_convergence}
\begin{tabular}{@{}lccccc@{}}
\toprule
Dropout rate & Final accuracy & Final F1 & Rounds to convergence & Stability gap & Prototype drift \\
\midrule
\tabrows{generated/dropout_convergence.tex}
\bottomrule
\end{tabular}
\end{table*}

% TODO(result): Fill this table after simulations are completed. Do not fabricate numbers.
\begin{table*}[t]
\centering
\caption{Stale-prototype robustness under bounded-staleness episodic updates. Single run (seed 0) per cell, including reference rows; journal protocol, 60 rounds, final-round test metrics.}
\label{tab:staleness_convergence}
\begin{tabular}{@{}lccccc@{}}
\toprule
Staleness & Final accuracy & Final F1 & Convergence rate & Prototype drift & Stability bound term \\
\midrule
\tabrows{generated/staleness_convergence.tex}
\bottomrule
\end{tabular}
\end{table*}

% TODO(figure): Replace placeholder with final plot after simulations are completed.
\begin{figure}[t]
    \centering
    \figorph{figures/fig_dropout_effect.pdf}{Placeholder: dropout-rate effect on convergence and drift (final plot pending simulations)}
    \caption{Effect of prototype dropout rate on convergence, stability, and prototype drift.}
    \label{fig:dropout_effect}
\end{figure}

% TODO(figure): Replace placeholder with final plot after simulations are completed.
\begin{figure}[t]
    \centering
    \figorph{figures/fig_staleness_effect.pdf}{Placeholder: staleness effect on convergence and drift (final plot pending simulations)}
    \caption{Effect of stale prototype updates on convergence, stability, and prototype drift.}
    \label{fig:staleness_effect}
\end{figure}

% TODO(result): Fill this table after simulations are completed. Do not fabricate numbers.
\begin{table*}[t]
\centering
\caption{Robustness under serverless-style event loss, timeout, cold-start, and partial-execution stresses. Single run (seed 0) per cell, including reference rows; journal protocol, 60 rounds, final-round test metrics.}
\label{tab:serverless_robustness}
\begin{tabular}{@{}lccccc@{}}
\toprule
Stress condition & Accuracy & F1-score & Prototype drift & Convergence rounds & Recovery gap \\
\midrule
\tabrows{generated/serverless_robustness.tex}
\bottomrule
\end{tabular}
\end{table*}

% TODO(result): Fill this table after simulations are completed. Do not fabricate numbers.
\begin{table*}[t]
\centering
\caption{IID and non-IID robustness for \hsfp and \ehsfp on classification datasets under the journal protocol (200 clients, 60 rounds, single run with seed 0 per cell, final-round test metrics); ``Non-IID mild'' and ``Non-IID severe'' are $dir(0.7)$ and $dir(0.3)$. Both rows use the corrected journal backbone; conference-reference \hsfp accuracies are reported only in Table~\ref{tab:main_classification}. Communication is in GB; n.c.\ marks runs that did not meet the convergence rule within 60 rounds.\PVpartialNote}
\label{tab:noniid_robustness}
% NOTE(provenance): Non-\pending values in this table are reproduced verbatim from the completed experiments of the accepted ECCV 2026 conference version (H-SFP paper); E-HSFP values pending journal simulation campaign.
\scriptsize
\setlength{\tabcolsep}{3pt}
\begin{tabular}{@{}lllccccc@{}}
\toprule
Dataset & Partition & Method & Accuracy & F1-score & Convergence rounds & Prototype drift & Communication \\
\midrule
\tabrows{generated/noniid_robustness.tex}
\bottomrule
\end{tabular}
\end{table*}

\subsection{Ablation Studies}
The ablation studies first reproduce the conference analysis of the H-SFP components in Table~\ref{tab:component_ablation}, then isolate the journal-only E-HSFP modules in Table~\ref{tab:ehsfp_component_ablation}. The conference-reference ablation shows that the complete H-SFP design improves CIFAR-100 IID accuracy by 19.63 points over HSFL without prototyping and by 2.75 points over the mean-only synthesis variant. The journal-only ablations are single runs with seed 0 on CIFAR-100 IID, so differences below roughly 0.6 points are within seed variance.

In the cumulative component ablation (Table~\ref{tab:ehsfp_component_ablation}, Figure~\ref{fig:ablation_modules}), adding episodic memory, prototype dropout, or reliability-aware aggregation alone leaves accuracy at or slightly below the \hsfp baseline (43.80--44.17 vs.\ 44.62\%). The first clear gain comes from replay consistency, which raises accuracy to 46.99\% ($+2.37$ points). The full \ehsfp configuration reaches 48.52\% ($+3.90$ points) with 4\% less communication (120.1 vs.\ 125.5~GB). The isolated residual-generator row was not run. In the synthesis ablation (Table~\ref{tab:synthesis_ablation}), low-rank covariance synthesis outperforms diagonal Gaussian synthesis by 1.38 points (48.52 vs.\ 47.14\%). This suggests that cross-feature dependencies, which the diagonal model of Section~\ref{sec:convergence} omits, matter in practice.

Memory capacity up to 10 has no measurable effect (Table~\ref{tab:memory_size_ablation}). Capacities 0, 1, 5, and 10 give 46.77--47.32\% accuracy, within 0.55 points, and capacity 20 gives \PVmemTwentyAcc. The additional memory overhead was not separately measured. The aggregation rules are likewise indistinguishable (Table~\ref{tab:aggregation_comparison}). Clean accuracy lies within 43.93--44.85\%. Under dropout and staleness, reliability-aware aggregation is highest (45.01 and 45.43\%), but only 0.29 and 0.09 points ahead of simple averaging, well within single-seed variance. Figure~\ref{fig:memory_effect} shows the memory effect over rounds, and Figure~\ref{fig:reliability_weights} shows the reliability-aware aggregation weights.

\begin{table*}[t]
\centering
\caption{Conference-reference H-SFP component ablation at $I_c=5$ and $I_e=10$ across 200 distributed devices. Results are accuracy (\%), mean $\pm$ standard deviation over five independent runs, reproduced verbatim from the accepted ECCV 2026 paper.}
\label{tab:component_ablation}
\begin{tabular}{@{}lcccc@{}}
\toprule
Variant & CIFAR-100 IID & CIFAR-100 $dir(0.7)$ & ImageNet-1K IID & ImageNet-1K $dir(0.7)$ \\
\midrule
HSFL, no prototyping & 35.47 {\scriptsize $\pm$ 1.35} & 8.25 {\scriptsize $\pm$ 1.15} & 24.6 {\scriptsize $\pm$ 0.9} & 14.8 {\scriptsize $\pm$ 1.1} \\
$+$ Contrastive loss & 41.20 {\scriptsize $\pm$ 1.25} & 8.95 {\scriptsize $\pm$ 1.05} & 25.4 {\scriptsize $\pm$ 0.8} & 16.3 {\scriptsize $\pm$ 1.0} \\
$+$ Synthetic data & 47.60 {\scriptsize $\pm$ 1.15} & 9.45 {\scriptsize $\pm$ 0.95} & 26.1 {\scriptsize $\pm$ 0.7} & 18.2 {\scriptsize $\pm$ 0.9} \\
Full, mean only & 52.35 {\scriptsize $\pm$ 1.05} & 9.85 {\scriptsize $\pm$ 0.90} & 26.8 {\scriptsize $\pm$ 0.5} & 20.5 {\scriptsize $\pm$ 0.8} \\
Full, purely generative & 48.20 {\scriptsize $\pm$ 1.10} & 9.15 {\scriptsize $\pm$ 1.00} & 25.8 {\scriptsize $\pm$ 0.7} & 19.4 {\scriptsize $\pm$ 0.9} \\
\textbf{Full H-SFP} & \textbf{55.10 {\scriptsize $\pm$ 0.95}} & \textbf{10.48 {\scriptsize $\pm$ 0.85}} & \textbf{27.9 {\scriptsize $\pm$ 0.3}} & \textbf{22.4 {\scriptsize $\pm$ 0.6}} \\
\bottomrule
\end{tabular}
\end{table*}

% TODO(result): Fill only after five-seed frozen-protocol runs are complete.
\begin{table*}[t]
\centering
\caption{Journal-only E-HSFP component ablation under the frozen episodic protocol. Every row will use identical partitions, seeds, round budgets, aggregation intervals, and model-selection rules. Single run (seed 0) per cell, including reference rows; journal protocol, 60 rounds, final-round test metrics.}
\label{tab:ehsfp_component_ablation}
\begin{tabular}{@{}lccccc@{}}
\toprule
Variant & Accuracy & Macro F1 & Communication & Convergence rounds & Stability \\
\midrule
\tabrows{generated/ehsfp_component_ablation.tex}
\bottomrule
\end{tabular}
\end{table*}

% INTERNAL DIAGNOSTIC NOTE (not manuscript evidence): single-seed 60-round
% CIFAR-100 results are 8.10/30414.23 MB (baseline), 11.28/1430.47 MB
% (memory), 11.20/1426.39 MB (memory+dropout), and 11.26/1430.47 MB
% (memory+dropout+reliability). The communication collapse is an unresolved
% accounting artifact. HAM10000 single-seed accuracies are 54.72, 50.77, and
% 55.82 for the latter three presets. ISIC proxy seed-20260714 IoU/Dice values
% are 50.28/64.95, 47.92/62.63, 47.83/62.37, and 46.74/61.69. None of these
% values may be promoted to the table above without the frozen five-seed rerun.

% TODO(result): Fill this table after simulations are completed. Do not fabricate numbers.
\begin{table*}[t]
\centering
\caption{Episodic prototype memory-size ablation under the dropout and staleness stress protocol. Single run (seed 0) per cell, including reference rows; journal protocol, 60 rounds, final-round test metrics.\PVpartialNote}
\label{tab:memory_size_ablation}
\begin{tabular}{@{}lccccc@{}}
\toprule
Memory size & Accuracy & F1-score & Prototype drift & Convergence rounds & Additional memory overhead \\
\midrule
\tabrows{generated/memory_size_ablation.tex}
\bottomrule
\end{tabular}
\end{table*}

% TODO(result): Fill this table after simulations are completed. Do not fabricate numbers.
\begin{table*}[t]
\centering
\caption{Aggregation-rule comparison under dropout and staleness stress protocols. Single run (seed 0) per cell, including reference rows; journal protocol, 60 rounds, final-round test metrics.}
\label{tab:aggregation_comparison}
\scriptsize
\setlength{\tabcolsep}{3pt}
\begin{tabular}{@{}lccccc@{}}
\toprule
Aggregation rule & Accuracy & F1-score & Robustness under dropout & Robustness under staleness & Communication \\
\midrule
\tabrows{generated/aggregation_comparison.tex}
\bottomrule
\end{tabular}
\end{table*}

% TODO(result): Fill this table after simulations are completed. Do not fabricate numbers.
\begin{table*}[t]
\centering
\caption{Prototype-synthesis ablation under the shared episodic training protocol. Single run (seed 0) per cell, including reference rows; journal protocol, 60 rounds, final-round test metrics. $^\star$Wall-clock time measured on the previous host (RTX 3080 Ti, up to four concurrent runs); not comparable with times measured on the current host (RTX 6000 Ada, twelve concurrent runs).}
\label{tab:synthesis_ablation}
\scriptsize
\setlength{\tabcolsep}{3pt}
\begin{tabular}{@{}lccccc@{}}
\toprule
Synthesis module & Accuracy & F1-score & Prototype fidelity & Communication overhead & Training time \\
\midrule
\tabrows{generated/synthesis_ablation.tex}
\bottomrule
\end{tabular}
\end{table*}

% TODO(figure): Replace placeholder with final plot after simulations are completed.
\begin{figure}[t]
    \centering
    \figorph{figures/fig_ablation_modules.pdf}{Placeholder: module ablation trends (final plot pending simulations)}
    \caption{Module ablation trends for accuracy, convergence rounds, and stability.}
    \label{fig:ablation_modules}
\end{figure}

% TODO(figure): Replace placeholder with final plot after simulations are completed.
\begin{figure}[t]
    \centering
    \figorph{figures/fig_reliability_weights.pdf}{Placeholder: reliability-aware aggregation weights (final plot pending simulations)}
    \caption{Reliability-aware prototype aggregation weights across support, variance, age, and reference-distance factors.}
    \label{fig:reliability_weights}
\end{figure}

% TODO(figure): Replace placeholder with final plot after simulations are completed.
\begin{figure}[t]
    \centering
    \figorph{figures/fig_memory_effect.pdf}{Placeholder: episodic prototype memory effect over rounds (final plot pending simulations)}
    \caption{Effect of episodic prototype memory on convergence, stability, and prototype drift over rounds.}
    \label{fig:memory_effect}
\end{figure}

\subsection{Efficiency and Scalability}
This subsection evaluates how performance changes with client population and how communication and memory are distributed across tiers. Table~\ref{tab:scalability} carries over the matching H-SFP CIFAR-100 IID accuracy values from the conference scalability experiment; unsupported client counts and journal-only metrics remain pending. Table~\ref{tab:comm_memory} disambiguates the interval-(A) and interval-(C) H-SFP communication totals reported in the conference paper. Under that common 200-round CIFAR-100 analysis, H-SFP(C) uses 10.94~GB: approximately $353\times$ less than FedAvg, $402\times$ less than HierFL, $20.8\times$ less than SplitFed, $10.3\times$ less than HSFL(A), and $1.58\times$ less than FedProto. Under the journal protocol, \hsfp accuracy on CIFAR-100 IID falls steadily as the fixed dataset is spread over more clients (Table~\ref{tab:scalability}): 62.86\% with 20 clients, 56.09\% with 50, 51.08\% with 100, 44.62\% with 200, and 36.54\% with 500. Over the same range, total communication grows from 27.2 to 275.8~GB. Peak client GPU memory saturates at 563~MB from 100 clients onward, and rounds to convergence stay between 30 and 39. Figure~\ref{fig:comm_acc_tradeoff} shows the communication--accuracy tradeoff, and Figure~\ref{fig:mem_perf_tradeoff} shows the memory--performance tradeoff under the journal protocol. The E-HSFP entry of Table~\ref{tab:comm_memory} remains pending because no 200-round conference-protocol \ehsfp run was performed.

% TODO(result): Fill this table after simulations are completed. Do not fabricate numbers.
\begin{table*}[t]
\centering
\caption{Scalability of \hsfp on CIFAR-100 IID under the journal protocol (60 rounds, single run with seed 0 per cell, final-round test metrics); the total dataset size is fixed while the number of clients varies. Conference-reference accuracies for 20--200 clients (67.80, 64.50, 60.20, 55.10) are not pooled with these values. $^\star$Wall-clock time measured on the previous host (RTX 3080 Ti, up to four concurrent runs); not comparable with times measured on the current host (RTX 6000 Ada, twelve concurrent runs).}
\label{tab:scalability}
% NOTE(provenance): Non-\pending values in this table are reproduced verbatim from the completed experiments of the accepted ECCV 2026 conference version (H-SFP paper); E-HSFP values pending journal simulation campaign.
\scriptsize
\setlength{\tabcolsep}{3pt}
\begin{tabular}{@{}lcccccc@{}}
\toprule
Clients & Accuracy & F1-score & Comm. & Client mem. & Time & Conv. rounds \\
\midrule
\tabrows{generated/scalability.tex}
\bottomrule
\end{tabular}
\end{table*}

\begin{table}[t]
\centering
\caption{Conference-reference total communication on CIFAR-100 over 200 rounds. Values are the literal totals stated in the accepted paper or displayed in its communication figure. Ratios use H-SFP(C), the most communication-efficient configuration, as reference.}
\label{tab:comm_memory}
\begin{tabular}{@{}lcc@{}}
\toprule
Method & Comm. (GB) $\downarrow$ & Relative cost \\
\midrule
FedAvg & $\sim$3,865 & $\sim$353.3$\times$ \\
HierFL & 4,392.44 & 401.5$\times$ \\
SplitFed & $\sim$228.0 & $\sim$20.8$\times$ \\
FedProto & 17.25 & 1.58$\times$ \\
HSFL (A) & 113.12 & 10.34$\times$ \\
H-SFP (A) & 13.8 & 1.26$\times$ \\
\textbf{H-SFP (C)} & \textbf{10.94} & \textbf{1.00$\times$} \\
E-HSFP & \pending & \pending \\
\bottomrule
\end{tabular}
\end{table}

% INTERNAL DIAGNOSTIC NOTE (not manuscript evidence): the unvalidated 60-round
% legacy-seed CIFAR-100 grid reports total_comm_MB of 222183.80 (SplitFL),
% 757374.46 (HSFL), 565462.98 (Federated), 115662.88 (HierFL), approximately
% 36700 (HeteroSFL, seed-dependent), and 756.16 (H-SFP). Although this preserves
% a large communication advantage for H-SFP, its corresponding accuracy is only
% 11.49+-0.25, so the grid cannot support a final communication--accuracy claim
% until the configuration mismatch and accounting checks are resolved.

% TODO(figure): Replace placeholder with final plot after simulations are completed.
\begin{figure}[t]
    \centering
    \figorph{figures/fig_comm_acc_tradeoff.pdf}{Placeholder: communication versus accuracy tradeoff (final plot pending simulations)}
    \caption{Communication versus accuracy tradeoff under the shared classification protocol.}
    \label{fig:comm_acc_tradeoff}
\end{figure}

% TODO(figure): Replace placeholder with final plot after simulations are completed.
\begin{figure}[t]
    \centering
    \figorph{figures/fig_mem_perf_tradeoff.pdf}{Placeholder: memory versus performance tradeoff (final plot pending simulations)}
    \caption{Memory versus performance tradeoff across client, edge, and cloud tiers.}
    \label{fig:mem_perf_tradeoff}
\end{figure}